\documentclass[tikz,border=8pt]{standalone}
\usepackage{ctex}
\usepackage{amsmath}
\usepackage{amssymb}
\usetikzlibrary{calc, arrows.meta, positioning}
\begin{document}
\begin{tikzpicture}[scale=1.0, thick,
  map/.style={-{Stealth[length=2.2mm]}, gray, thin},
  stem/.style={-{Stealth[length=2.2mm]}, blue!70!black, ultra thick},
]
  % ===== 标题 =====
  \node[font=\bfseries] at (4.7,5.25)
    {随机变量 $X$：把样本空间的概率「搬运」到数轴 $\to$ 分布 $P_X$};

  % ===== 左：样本空间 Omega（抛两枚硬币，X = 正面数）=====
  \draw[black, thick] (0,1.9) ellipse (1.55 and 2.25);
  \node[font=\small\bfseries, above] at (0,4.15) {$\Omega$（抛两枚硬币）};

  % 四个等概率结果
  \foreach \y/\name in {3.4/正正, 2.4/正反, 1.4/反正, 0.4/反反} {
    \fill[black] (0,\y) circle (1.8pt);
    \node[font=\small, left=2pt] at (-0.05,\y) {\name};
    \node[font=\scriptsize, gray, right=2pt] at (0.05,\y) {$\tfrac14$};
  }

  % ===== 右：数轴 R 与诱导分布 P_X 的概率「砝码」=====
  \draw[->, thick] (5,1.4) -- (9.6,1.4) node[right, font=\small] {$\mathbb{R}$};
  \foreach \x/\lbl in {5.9/0, 7.25/1, 8.6/2} {
    \draw (\x,1.32) -- (\x,1.48);
    \node[font=\small, below] at (\x,1.28) {\lbl};
  }
  % P_X 概率质量（竖直「砝码」）：1/4 -> 0.7 单位
  \draw[stem] (5.9,1.4) -- (5.9,2.1);   \node[font=\scriptsize, blue!70!black, above] at (5.9,2.1) {$\tfrac14$};
  \draw[stem] (7.25,1.4) -- (7.25,2.8);  \node[font=\scriptsize, blue!70!black, above] at (7.25,2.8) {$\tfrac12$};
  \draw[stem] (8.6,1.4) -- (8.6,2.1);   \node[font=\scriptsize, blue!70!black, above] at (8.6,2.1) {$\tfrac14$};
  \node[font=\small\bfseries, blue!70!black, anchor=west] at (8.8,2.6) {$P_X$};

  % ===== 搬运箭头：每个 omega -> 它的像 X(omega) =====
  \draw[map] (0.4,0.4)  to[bend right=8]  (5.9,1.36);   % 反反 -> 0
  \draw[map] (0.45,2.4) to[bend left=4]   (7.25,1.36);  % 正反 -> 1
  \draw[map] (0.45,1.4) to[bend right=4]  (7.25,1.36);  % 反正 -> 1
  \draw[map] (0.4,3.4)  to[bend left=10]  (8.6,1.36);   % 正正 -> 2

  % ===== 半透明注释框（置于数轴下方空白处，避开砝码峰值）=====
  \node[font=\scriptsize, fill=yellow!18, draw=yellow!60!black, rounded corners=2pt,
        align=center, fill opacity=0.85, text opacity=1, anchor=north]
    at (7.25,0.95)
    {$\displaystyle P_X(B)=P(X\in B)=P\big(X^{-1}(B)\big)$\\[1pt]「正反」与「反正」两份概率在 $X=1$ 处合并相加};

\end{tikzpicture}
\end{document}
